%%%%%%%%%%%%%%%%%%%%%%%%%%%%%%%%%%
% Appendix
%%%%%%%%%%%%%%%%%%%%%%%%%%%%%%%%%%

\section{Proof of theorem~\ref{thm:approx_expected_precision}}
\label{sec:proof_approx_expected_precision}
Given $p$ positives and $n$ negatives, by \cref{thm:approx_expected_precision} an approximate set with a false positive rate $\fprate$ and a false negative rate $\fnrate$ has an \emph{expected} precision given \emph{approximately} by
\begin{equation*}
    \ppv(\fnrate, \fprate ; n, p) \approx \frac{\overline{t}}{\overline{t} + \overline{f}} +
    \frac{\overline{t} \sigma_{\!f}^2 - \overline{f} \sigma_{\!t}^2}{\left(\overline{t} + \overline{f}\right)^3},
    \tag{\ref{eq:approx_expected_precision} revisited}
\end{equation*}
where $\overline{t} = p \tprate$ is the \emph{expected} number of \emph{true positives}, $\overline{f} =  n \fprate$ is the \emph{expected} number of \emph{false positives}, $\sigma_{\!t}^2 = p \fnrate \tprate$ is the variance of the number of \emph{true positives}, and $\sigma_{\!f}^2 = n \fprate \tnrate$ is the variance of the number of \emph{false positives}.
\begin{proof}
The positive predictive value is a random variable given by
\begin{equation}
    \frac{\TP_p}{\TP_p + \FP_n}.
\end{equation}
We are interested in the \emph{expected} positive predictive value,
\begin{equation}
    \ppv(\fprate,\tprate) = \Expect{\frac{\TP_p}{\TP_p + \FP_n}}.
\end{equation}
This expectation is of a non-linear function of random variables, which is problematic so we choose to approximate the expectation.

Let the \emph{positive predictive value} function be denoted by
\begin{equation}
    f(t, f) = \frac{t}{t + f},
\end{equation}
where $t$ is the number of true positives and $f$ is the number of false positives.
We approximate this function with a second-order Taylor series.
The gradient of $f$ is given by
\begin{equation}
    \nabla f(t,f) =
    \frac{1}{(t + f)^2}
    \begin{bmatrix}
        f\\
        -t\\
    \end{bmatrix}
\end{equation}
and the Hessian of $f$ is given by
\begin{equation}
    \mathcal{H}(t,f) =
    \frac{1}{(t + f)^3}
    \begin{bmatrix}
        -2 f & t-f\\
        t-f & 2 t \\
    \end{bmatrix}
\end{equation}

A second-order Taylor approximation $\Fun{g}$ of $f$ about the expected value of $\TP_p$, denoted by $\overline{t}$, and the expected value of $\FP_n$, denoted by $\overline{f}$, is given by
\begin{equation}
    \Fun{g}(t,f) =
    f\left(\overline{t},\overline{f}\right) + \nabla f(\overline{t},\overline{f})^{\intercal}
    \begin{bmatrix}
        t - \overline{t}\\
        f - \overline{f}\\
    \end{bmatrix}
    + \frac{1}{2}
    \begin{bmatrix}
        t - \overline{t}\\
        f - \overline{f}\\
    \end{bmatrix}^{\intercal}
    \mathcal{H}(\overline{t},\overline{f})
    \begin{bmatrix}
        t - \overline{t}\\
        f - \overline{f}\\
    \end{bmatrix}
\end{equation}
As a function of random variables $\TP_p$ and $\FP_n$, $\Fun{g}\!\left(\TP_p,\FP_n\right)$ is a random variable.
Since $\Expect{\TP_p - \overline{t}} = 0$ and $\Expect{\FP_n - \overline{f}} = 0$, we immediately simplify the expectation of $\Fun{g}$ to
\begin{equation}
\label{eq:proof_hess1}
    \Expect{\Fun{g}(\TP_p,\FP_n)} = \frac{\overline{t}}{\overline{t}+\overline{f}} + \frac{\Expect{A}}{(\overline{t} + \overline{f})^3}
\end{equation}
where
\begin{equation}
    A = \frac{1}{2}
    \begin{bmatrix}
        \TP_p - \overline{t}\\
        \FP_n - \overline{f}
    \end{bmatrix}^{\intercal}
    \begin{bmatrix}
        -2 \overline{f} \left(\TP_p - \overline{t}\right) + \left(\overline{t}-\overline{f}\right)\left(\FP_n - \overline{f}\right)\\
        \left(\overline{t}-\overline{f}\right)\left(\TP_p - \overline{t}\right) + 2 \overline{t}\left(\FP_n - \overline{f}\right)
    \end{bmatrix}
\end{equation}
Expanding the quadratic form yields
\begin{equation}
    A = -\overline{f} \left(\TP_p - \overline{t}\right)^2 +
    \left(\overline{t}-\overline{f}\right)\left(\TP_p -
    \overline{t}\right)\left(\FP_n - \overline{f}\right) +
    \overline{t}\left(\FP_n - \overline{f}\right)^2
\end{equation}
As a linear operator, the expectation of $A$ is equivalent to
\begin{equation}
    \Expect{A} = -\overline{f} \Expect{(\TP_p - \overline{t})^2} + \left(\overline{t}-\overline{f}\right)\Expect{\left(\FP_n - \overline{f}\right)\left(\TP_p - \overline{t}\right)} + \overline{t}\Expect{(\FP_n - \overline{f})^2}
\end{equation}
By definition, $\Expect{(\TP_p - \overline{t})^2}$ is the variance of $\TP_p$, $\Expect{(\FP_n - \overline{f})^2}$ is the variance of $\FP_n$, and $\Expect{\left(\FP_n - \overline{f}\right)\left(\TP_p - \overline{t}\right)}$ is the covariance of $\TP_p$ and $\FP_n$, which is $0$ since they are independent.
Making these substitutions results in
\begin{equation}
    \Expect{A} = \overline{t} \Var{\FP_n} - \overline{f} \Var{\TP_p}
\end{equation}
Substituting this result into \cref{eq:proof_hess1} yields
\begin{equation}
    \Expect{\Fun{g}(\TP_p,\FP_n)} =
    \frac{\overline{t}}{\overline{t}+\overline{f}} +
    \frac{-\overline{f} \Var{\TP_p} +
    \overline{t}\Var{\FP_n}}{(\overline{t} + \overline{f})^3}
\end{equation}
By \cref{thm:fpbinom}, $\FP_n$ is binomially distributed with a mean $n \fprate$ and a variance $n \fprate \tnrate$ and by \cref{cor:tpbinom}, $\TP_p$ is binomially distributed with a mean $p \tprate$ and a variance $p \fnrate \tprate$.
Substituting $\Var{\FP_n} = n \fprate(1-\fprate) = \overline{f}_p(1-\fprate) = \sigma_{f_p}^2$ and $\Var{\TP_p} = p \tprate(1-\tprate) = \overline{t}_p(1-\tprate) = \sigma_{t_p}^2$ completes the proof.
\end{proof}

\section{Proof of theorem~\ref{thm:f1_approx}}
\label{sec:proof_f1}
The proof follows the same structure as the PPV proof (\cref{sec:proof_approx_expected_precision}) with the substitution $h(t,f) = 2t/(t + f + \p)$ in place of $f(t,f) = t/(t+f)$.
\begin{proof}
The $\fone$-score is a random variable $\Fone = h(\TP_\p, \FP_\n)$ where
\begin{equation}
    h(t, f) = \frac{2t}{t + f + \p}.
\end{equation}
The gradient of $h$ is
\begin{equation}
    \nabla h(t,f) =
    \frac{2}{(t + f + \p)^2}
    \begin{bmatrix}
        f + \p\\
        -t\\
    \end{bmatrix}
\end{equation}
and the Hessian is
\begin{equation}
    \mathcal{H}(t,f) =
    \frac{2}{(t + f + \p)^3}
    \begin{bmatrix}
        -2(f+\p) & t - f - \p\\
        t - f - \p & 2 t \\
    \end{bmatrix}.
\end{equation}
The second-order Taylor expansion about the means $\overline{t} = \p\tprate$, $\overline{f} = \n\fprate$, letting $s = \overline{t} + \overline{f} + \p$, gives
\begin{equation}
    \Expect{h(\TP_\p, \FP_\n)} \approx h(\overline{t}, \overline{f}) + \frac{1}{2}\operatorname{tr}\!\bigl(\mathcal{H}(\overline{t},\overline{f})\,\Sigma\bigr)
\end{equation}
where $\Sigma = \operatorname{diag}(\sigma_t^2, \sigma_f^2)$ and the cross terms vanish by independence.
Expanding:
\begin{equation}
    \frac{1}{2}\operatorname{tr}\!\bigl(\mathcal{H}\,\Sigma\bigr) =
    \frac{1}{2}\Bigl(\frac{-4(\overline{f}+\p)}{s^3}\,\sigma_t^2 +
    \frac{4\overline{t}}{s^3}\,\sigma_f^2\Bigr) =
    \frac{2\overline{t}\,\sigma_f^2 - 2(\overline{f}+\p)\sigma_t^2}{s^3}.
\end{equation}
Adding $h(\overline{t},\overline{f}) = 2\overline{t}/s$ completes the expected-value formula.

For the variance, the first-order delta method gives $\Var{\Fone} \approx (\nabla h)^\intercal \Sigma\,(\nabla h)$:
\begin{equation}
    \Var{\Fone} \approx \frac{4(\overline{f}+\p)^2 \sigma_t^2 + 4\overline{t}^{\,2}\,\sigma_f^2}{s^4}.
\end{equation}
Substituting $\sigma_t^2 = \p\tprate(1-\tprate)$ and $\sigma_f^2 = \n\fprate(1-\fprate)$ yields \cref{eq:f1_approx,eq:var_f1}.
\end{proof}

\section{Sampling distribution of arbitrary functions}
\label{app:samp}
Suppose we have a function $f : \PS{\Set{X}[1]} \times \cdots \times \PS{\Set{X}[q]} \to \Set{Y}$, and we are interested in evaluating the loss when we replace one or more of the objective input sets with particular corresponding random approximate sets.
The result of this substitution, as previously described, is a probability distribution over $\Set{Y}$.

The probability distribution of random approximate sets are precisely given; therefore, we may estimate the distribution of any function of random approximate sets by generating the random approximate sets and applying the function of interest.

Consider the $m$-by-$q$ matrix where the $(i,j)$-th element is the random approximate set $\tilde{A}_{i\,j}$ such that each random element of the matrix is independently and each column is identically distributed.
If we apply $\Fun{g}$ to each row of the matrix,
\begin{equation}
\RV{Y}_i = \Fun{g}\left(\tilde{A}_{i\,1},\ldots,\tilde{A}_{i\,q}\right)
\end{equation}
for $i=1,\ldots,m$, we generate $m$ i.i.d. random elements $\RV{Y_1},\ldots,\RV{Y_m}$.

If $\Set{Y}$ is a measure space (discrete or continuous), consider the random variable
\begin{equation}
\RV{\overline{Y}_m} = \frac{1}{m} \sum_{i=1}^m \RV{Y_i}.
\end{equation}
If $\RV{Y_1}$ has a well-defined mean $\mu = \Expect{\RV{Y_1}}$ and variance $\sigma^2 = \Var{\RV{Y_1}}$, then by the central limit theorem
\begin{equation}
\sqrt{m}\left(\RV{\overline{Y}_m} - \mu\right) \xrightarrow{d} \normdist\!\left(0, \sigma^2\right)
\end{equation}
as $m \to \infty$.
That is, for large $m$ the sample mean $\RV{\overline{Y}_m}$ is approximately normally distributed with mean $\mu$ and variance $\sigma^2 / m$.

This framework underlies the Monte Carlo validation of the delta method approximations in \cref{fig:ppv_theory_vs_sim,fig:f1_theory_vs_sim}: for each parameter value, $m = 50\,000$ independent random approximate sets are generated, and the sample mean of the measure is compared to the theoretical prediction.
The large sample size ensures that the Monte Carlo standard error is negligible relative to the quantities being estimated.
